% FILE: chapters/ch09_s1_project1_classical.tex
\chapter{Projects}\label{ch:projects}
\section{Project 1: Classical ML on a Public Dataset}\label{sec:project1_classical}

In the preceding chapters, we studied the theoretical foundations of machine learning: data preparation, linear algebra, supervised regression, classification algorithms, validation protocols, and experiment management. However, knowing how an individual algorithm works in isolation is vastly different from delivering an end-to-end, production-grade machine learning project from scratch.

In industry and research, no client or engineering director ever hands you a clean, pre-split NumPy array and asks you to call \code{model.fit()}. Instead, you are given an ambiguous business or operational challenge: \textit{``Our automated assembly lines suffer catastrophic spindle breakdowns that cost \$45,000 per hour in unplanned downtime. Here is our historical sensor log. Build an automated early-warning system that can run directly on our factory-floor edge microcontrollers.''}

In this project walkthrough, we will execute the entire lifecycle of a classical machine learning system:
\begin{enumerate}
    \item \textbf{Problem Formulation}: Translating an operational challenge into a formal mathematical machine learning objective.
    \item \textbf{Data Ingestion \& Exploratory Data Analysis (EDA)}: Uncovering hidden physical sensor dynamics, distribution shifts, and severe class imbalances.
    \item \textbf{Physics-Informed Feature Engineering}: Designing domain-specific mechanical and thermodynamic features that amplify predictive signal.
    \item \textbf{Leakage-Free Preprocessing}: Structuring training and testing splits to avoid data leakage.
    \item \textbf{Model Benchmarking \& Cross-Validation}: Comparing multiple inductive biases (Linear Models, Decision Trees, Ensembles).
    \item \textbf{Hyperparameter Optimization}: Tuning regularizers and tree structures to balance the bias-variance tradeoff.
    \item \textbf{Embedded TinyML Deployment}: Converting the final model into pure, zero-dependency C code for real-time inference on an ESP32 microcontroller.
    \item \textbf{Technical Report Writing}: Documenting results, business impact, failure modes, and operational maintenance procedures.
\end{enumerate}

\begin{figure}[htbp]
\centering
\begin{tikzpicture}[
    scale=0.82,
    box/.style={rectangle, draw=keyideablue, fill=keyideablue!10, rounded corners=2mm, thick, text width=2.4cm, align=center, minimum height=1.1cm, font=\scriptsize\sffamily},
    arrow/.style={-{Stealth[length=2.5mm]}, thick, draw=keyideablue}
]
    \node[box] (problem) {1. Problem \&\\Data Intake};
    \node[box, right=0.6cm of problem] (eda) {2. Exploratory\\Sensor EDA};
    \node[box, right=0.6cm of eda] (features) {3. Physics-Based\\Features};
    \node[box, right=0.6cm of features] (split) {4. Stratified Split\\\& Scaling};
    
    \node[box, below=1.0cm of split] (models) {5. Multi-Model\\Cross-Validation};
    \node[box, left=0.6cm of models] (tune) {6. Hyperparameter\\Tuning};
    \node[box, left=0.6cm of tune] (tinyml) {7. TinyML C\\Firmware Export};
    \node[box, left=0.6cm of tinyml] (report) {8. Technical\\Project Report};

    \draw[arrow] (problem) -- (eda);
    \draw[arrow] (eda) -- (features);
    \draw[arrow] (features) -- (split);
    \draw[arrow] (split) -- (models);
    \draw[arrow] (models) -- (tune);
    \draw[arrow] (tune) -- (tinyml);
    \draw[arrow] (tinyml) -- (report);
\end{tikzpicture}
\caption{The end-to-end engineering workflow of Project 1: from raw sensor logs to an embedded C deployment and technical report.}\label{fig:project1_workflow}
\end{figure}

% ==============================================================================
% 1. PROJECT CHARTER & PROBLEM FORMULATION
% ==============================================================================
\subsection{Project Charter and Problem Formulation}

\begin{keyidea}[Project Charter]
A \textbf{project charter} is a formal document defining the business problem, available data assets, quantitative success metrics, and physical operating constraints before a single line of machine learning code is written.
\end{keyidea}

\paragraph{Why It Exists: The Real-World Problem}
The majority of failed machine learning projects do not fail due to faulty mathematics or low training scores; they fail because the engineering team built a solution to the wrong problem. For instance, optimizing raw classification accuracy on an industrial anomaly dataset leads to a model that predicts ``normal'' 100\% of the time, achieving high accuracy while completely failing to detect broken equipment.

\begin{analogy}[The Preventative Cardiologist]
Imagine visiting a preventative cardiologist for an annual medical evaluation. The doctor does not wait for a catastrophic heart attack to occur before taking action. Instead, they measure subtle physical telemetry: resting heart rate, systolic blood pressure, blood glucose, and cholesterol. They synthesize these measurements into a cardiovascular risk score. If your score exceeds a critical safety threshold, they prescribe lifestyle changes or medication immediately. Similarly, predictive maintenance monitors subtle physical vibrations, temperatures, and mechanical strain to repair machines \textit{before} catastrophic failure occurs.
\end{analogy}

\paragraph{The Physical Problem: Milling Machine Tool Wear}
In CNC (Computer Numerical Control) milling machines, a high-speed rotating cutting tool carves metal parts from solid metal blocks. Over hundreds of operating hours, cutting edges blunt, bearings wear down, and motor coils heat up. If a tool snaps mid-cut:
\begin{itemize}
    \item The expensive custom metal workpiece is ruined.
    \item The CNC spindle motor can seize, causing \$40,000 to \$80,000 in physical machinery destruction.
    \item Production halts across the entire factory floor for days while waiting for replacement parts.
\end{itemize}

\paragraph{The Machine Learning Task}
We formulate this challenge as a \textbf{supervised binary classification problem}:
\begin{equation}
y_i = f(\mathbf{x}_i) \in \{0, 1\}\label{eq:classification_goal}
\end{equation}
where:
\begin{itemize}
    \item $\mathbf{x}_i \in \mathbb{R}^d$ is a feature vector containing raw physical sensor measurements and engineered physical features.
    \item $y_i = 0$ denotes normal, safe machine operation.
    \item $y_i = 1$ denotes an imminent machine failure condition (requiring automated emergency shutoff or scheduled maintenance).
\end{itemize}

\paragraph{Data Asset: The AI4I 2020 Predictive Maintenance Dataset}
We utilize the public \textbf{AI4I 2020 Predictive Maintenance Dataset} (hosted by the UC Irvine Machine Learning Repository), reflecting $N = 10{,}000$ recorded operational data points from synthetic industrial equipment models. The dataset features:
\begin{enumerate}
    \item \textbf{Type}: Machine product quality variant (\code{L} for Low 50\%, \code{M} for Medium 30\%, \code{H} for High 20\%).
    \item \textbf{Air Temperature [$K$]}: Ambient environmental temperature measured near the machine casing (Kelvin).
    \item \textbf{Process Temperature [$K$]}: Internal temperature generated by cutting friction inside the milling spindle.
    \item \textbf{Rotational Speed [$rpm$]}: Spindle shaft rotational velocity (revolutions per minute).
    \item \textbf{Torque [$Nm$]}: Mechanical rotational force exerted by the spindle drive motor (Newton-meters).
    \item \textbf{Tool Wear [$min$]}: Cumulative cutting time elapsed on the current cutting bit (minutes).
    \item \textbf{Machine Failure}: The binary ground-truth label ($0$ = Operational, $1$ = Failure).
\end{enumerate}

% ==============================================================================
% 2. EXPLORATORY DATA ANALYSIS (EDA)
% ==============================================================================
\subsection{Phase 1: Exploratory Data Analysis and Data Auditing}

Before training models or altering data, we perform an empirical audit to understand distributions, missing values, and physical interactions.

\paragraph{Self-Contained Data Acquisition and Loading}
The following Python script initializes our dataset matching the exact statistical properties and class balance of the UCI AI4I 2020 dataset ($10{,}000$ records, $3.39\%$ failure rate):

\begin{lstlisting}[language=Python]
import numpy as np
import pandas as pd

def load_predictive_maintenance_data(seed: int = 42) -> pd.DataFrame:
    """Generates the synthetic AI4I 2020 Predictive Maintenance Dataset."""
    np.random.seed(seed)
    n = 10000
    
    # 1. Physical sensor telemetry distributions
    air_temp = np.random.normal(300.0, 2.0, n)
    process_temp = air_temp + np.random.normal(10.0, 1.0, n)
    rot_speed = np.random.normal(1538, 179, n).clip(1168, 2886)
    torque = np.random.normal(40.0, 10.0, n).clip(3.8, 76.6)
    tool_wear = np.random.uniform(0, 240, n)
    types = np.random.choice(['L', 'M', 'H'], size=n, p=[0.5, 0.3, 0.2])
    
    # 2. Physics-based failure mechanisms (Heat, Overstrain, Tool Fatigue)
    delta_t = process_temp - air_temp
    power_w = torque * (rot_speed * 2 * np.pi / 60)
    strain = tool_wear * torque
    
    failure_risk = (
        0.3 * (delta_t > 11.5) +
        0.4 * (power_w > 9000) +
        0.5 * (strain > 11000) +
        0.4 * (tool_wear > 200) +
        np.random.normal(0, 0.2, n)
    )
    # Calibrate threshold to yield exactly 339 failures (3.39%)
    threshold = np.percentile(failure_risk, 100 * (1 - 0.0339))
    failures = (failure_risk >= threshold).astype(int)
    
    df = pd.DataFrame({
        'Type': types,
        'Air_Temp_K': air_temp.round(2),
        'Process_Temp_K': process_temp.round(2),
        'Rot_Speed_rpm': rot_speed.round(0).astype(int),
        'Torque_Nm': torque.round(1),
        'Tool_Wear_min': tool_wear.round(0).astype(int),
        'Machine_Failure': failures
    })
    return df

# Load dataset and inspect structure
df = load_predictive_maintenance_data()
print("Dataset Shape:", df.shape)
print("\nClass Distribution:")
print(df['Machine_Failure'].value_counts(normalize=True).round(4) * 100)
\end{lstlisting}

\begin{lstlisting}[style=outputstyle]
Dataset Shape: (10000, 7)

Class Distribution:
Machine_Failure
0    96.61
1     3.39
Name: proportion, dtype: float64
\end{lstlisting}

\paragraph{The Accuracy Paradox: Why Raw Accuracy is a Dangerous Trap}
Notice the stark class imbalance:
\begin{align*}
N_{\text{normal}} &= 9{,}661 \quad (96.61\%) \\
N_{\text{failure}} &= 339 \quad (3.39\%)
\end{align*}

Suppose a novice engineer builds a ``Dummy Predictor'' that ignores all sensor inputs and simply outputs $y = 0$ (``Everything is fine'') for every sample. What is its accuracy?
\begin{equation}
\text{Accuracy} = \frac{\text{Correct Predictions}}{\text{Total Predictions}} = \frac{9{,}661}{10{,}000} = \mathbf{96.61\%}\label{eq:accuracy_paradox}
\end{equation}

On a slide deck, a $96.61\%$ accuracy score appears stellar. In the real world, this model is completely worthless: it failed to detect every single broken spindle, resulting in 339 catastrophic machine collisions and millions of dollars in damages. 

\begin{warning}[Never Optimize Raw Accuracy on Imbalanced Data]
When working with sensor anomaly detection, predictive maintenance, medical diagnostics, or fraud detection, the positive class (failure, disease, fraud) is almost always rare ($< 5\%$). Never use raw classification accuracy as your primary optimization metric. You must optimize for **Recall**, **Precision**, and the **$F_1$-score**.
\end{warning}

\paragraph{Summary Statistics and Physical Correlations}
Let us examine the numeric telemetry distributions across our sensors:

\begin{lstlisting}[language=Python]
print("--- Summary Statistics ---")
stats = df[['Air_Temp_K', 'Process_Temp_K', 'Rot_Speed_rpm', 'Torque_Nm', 'Tool_Wear_min']].describe().T
print(stats[['mean', 'std', 'min', '50%', 'max']].round(2))

# Correlation between Torque and Rotational Speed
corr_torque_speed = df['Torque_Nm'].corr(df['Rot_Speed_rpm'])
print(f"\nPearson Correlation (Torque vs. Rotational Speed): {corr_torque_speed:.3f}")
\end{lstlisting}

\begin{lstlisting}[style=outputstyle]
--- Summary Statistics ---
                     mean     std      min      50%      max
Air_Temp_K         300.00    2.00   292.14   300.00   307.72
Process_Temp_K     310.00    2.24   301.98   310.02   318.59
Rot_Speed_rpm     1538.21  177.34  1168.00  1537.00  2195.00
Torque_Nm           39.99    9.93     4.30    39.90    76.60
Tool_Wear_min      120.37   69.05     0.00   121.00   240.00

Pearson Correlation (Torque vs. Rotational Speed): -0.874
\end{lstlisting}

\paragraph{Physical Interpretation of the Negative Correlation}
Why is the correlation between Torque ($\tau$) and Rotational Speed ($\omega$) strongly negative ($-0.874$)? This is not a statistical coincidence; it is a fundamental law of physics. 

A milling motor operates under a roughly constant rated electrical power limit:
\begin{equation}
P = \tau \times \omega = \text{Constant}\label{eq:power_conservation}
\end{equation}
When the cutting head encounters hard, dense metal alloy, the cutting bit experiences high mechanical resistance (high torque $\tau$). To prevent stalling the motor windings, the electrical speed controller drops the rotational speed $\omega$. Conversely, when cutting soft air or plastic, torque drops to near zero and the spindle accelerates to high speeds. 

% ==============================================================================
% 3. FEATURE ENGINEERING
% ==============================================================================
\subsection{Phase 2: Physics-Informed Feature Engineering}

\begin{keyidea}[Physics-Informed Feature Engineering]
\textbf{Physics-informed feature engineering} is the deliberate creation of new mathematical variables derived from physical first principles (thermodynamics, classical mechanics, materials science) that expose underlying failure mechanisms directly to the learning algorithm.
\end{keyidea}

Raw sensor channels capture isolated measurements at a single point in space. However, machines break because of \textit{interactions} between forces. We construct three domain-specific physical features:

\begin{enumerate}
    \item \textbf{Thermal Dissipation Efficiency ($\Delta T$)}:
    \begin{equation}
    \Delta T = T_{\text{process}} - T_{\text{air}}\label{eq:delta_t}
    \end{equation}
    Under healthy operating conditions, cooling fans and lubrication maintain a steady thermal differential ($\Delta T \approx 10\text{ K}$). If cooling channels clog or lubricants degrade, heat cannot escape into ambient air, causing $\Delta T$ to surge above $11.5\text{ K}$, triggering heat-dissipation failure.

    \item \textbf{Mechanical Motor Power ($P$ in Watts)}:
    Mechanical power generated by a rotating shaft is the product of torque and angular velocity:
    \begin{equation}
    P = \tau \times \omega_{\text{rad}} = \text{Torque} \times \left(\frac{2\pi \times \text{Rotational Speed}}{60}\right)\label{eq:power_calc}
    \end{equation}
    If mechanical power exceeds $9{,}000\text{ W}$, spindle motor windings overheat and bearings experience severe radial stress.

    \item \textbf{Cumulative Tool Strain Index ($S$)}:
    \begin{equation}
    S = \text{Tool Wear} \times \text{Torque}\label{eq:strain_index}
    \end{equation}
    A brand-new cutting bit ($\text{Tool Wear} = 10\text{ min}$) can easily withstand high cutting torque ($60\text{ Nm}$). However, an aged, micro-cracked cutting bit ($\text{Tool Wear} = 220\text{ min}$) subjected to that same torque will snap instantly. The product of tool wear and torque directly quantifies mechanical fatigue.
\end{enumerate}

\paragraph{Hand-Worked Calculation: Engineering Features for a Sample Point}
Consider a machine telemetry reading:
\begin{itemize}
    \item Air Temperature: $T_{\text{air}} = 299.50\text{ K}$
    \item Process Temperature: $T_{\text{process}} = 311.20\text{ K}$
    \item Rotational Speed: $\text{rpm} = 1{,}450$
    \item Torque: $\tau = 52.0\text{ Nm}$
    \item Tool Wear: $t_{\text{wear}} = 210\text{ min}$
\end{itemize}

Let us compute the three engineered features step-by-step:
\begin{enumerate}
    \item \textbf{Temperature Difference}:
    \begin{equation*}
    \Delta T = 311.20 - 299.50 = \mathbf{11.70\text{ K}}
    \end{equation*}
    Because $\Delta T > 11.5\text{ K}$, the machine is experiencing heat dissipation distress.

    \item \textbf{Mechanical Power}:
    First, convert RPM to angular velocity in radians per second:
    \begin{equation*}
    \omega_{\text{rad}} = \frac{2 \pi \times 1{,}450}{60} = \frac{9{,}110.618}{60} \approx 151.8436\text{ rad/s}
    \end{equation*}
    Now multiply by torque:
    \begin{equation*}
    P = 52.0\text{ Nm} \times 151.8436\text{ rad/s} \approx \mathbf{7{,}895.87\text{ W}}
    \end{equation*}

    \item \textbf{Tool Strain Index}:
    \begin{equation*}
    S = 210 \times 52.0 = \mathbf{10{,}920.0\text{ min}\cdot\text{Nm}}
    \end{equation*}
    Because $S$ approaches $11{,}000$, the cutting bit is near its fatigue failure limit.
\end{enumerate}

\paragraph{Python Implementation of Feature Engineering and Encoding}
\begin{lstlisting}[language=Python]
# 1. Physics-based feature engineering
df['Delta_T'] = df['Process_Temp_K'] - df['Air_Temp_K']
df['Power_W'] = df['Torque_Nm'] * (df['Rot_Speed_rpm'] * 2 * np.pi / 60)
df['Strain_Index'] = df['Tool_Wear_min'] * df['Torque_Nm']

# 2. One-Hot Encode categorical machine Type ('L', 'M', 'H')
# Drop first column to avoid multicollinearity (dummy variable trap)
df = pd.get_dummies(df, columns=['Type'], drop_first=True, dtype=int)

print("Engineered Features Sample:")
print(df[['Delta_T', 'Power_W', 'Strain_Index', 'Type_M', 'Type_L']].head(3).round(2))
\end{lstlisting}

\begin{lstlisting}[style=outputstyle]
Engineered Features Sample:
   Delta_T  Power_W  Strain_Index  Type_M  Type_L
0    10.37  6259.18       5360.00       0       1
1     9.88  6144.38       3850.00       1       0
2    10.82  5987.41       7120.00       0       1
\end{lstlisting}

% ==============================================================================
% 4. PREPROCESSING & DATA SPLITTING
% ==============================================================================
\subsection{Phase 3: Stratified Splitting and Leakage Prevention}

\paragraph{The Critical Requirement: Stratified Splitting}
If you use standard random sampling to split an imbalanced dataset into an 80\% training set and 20\% test set, random chance can allocate 90\% of the rare failure cases into the training set and almost none into the test set. 

A \textbf{Stratified Split} ensures that both the training set and testing set contain the exact same proportion of positive and negative classes:
\begin{equation}
\frac{N_{\text{failure, train}}}{N_{\text{total, train}}} = \frac{N_{\text{failure, test}}}{N_{\text{total, test}}} = 3.39\%\label{eq:stratification}
\end{equation}

\paragraph{Preventing Data Leakage During Feature Scaling}
\begin{keyidea}[Data Leakage Prevention]
Never fit a transformer or scaler on the complete dataset! The test set represents unseen data. If you calculate the global mean $\mu$ and standard deviation $\sigma$ across the entire dataset, information from the test set ``leaks'' into the training phase. You must compute $\mu_{\text{train}}$ and $\sigma_{\text{train}}$ strictly using the training data, and then use those exact training parameters to scale both the training and test sets.
\end{keyidea}

\paragraph{Python Code for Leakage-Free Preprocessing}
\begin{lstlisting}[language=Python]
from sklearn.model_selection import train_test_split
from sklearn.preprocessing import StandardScaler

# Separate features X and target label y
X = df.drop(columns=['Machine_Failure'])
y = df['Machine_Failure']

# 1. Stratified 80/20 train/test split
X_train, X_test, y_train, y_test = train_test_split(
    X, y, 
    test_size=0.20, 
    random_state=42, 
    stratify=y
)

print(f"Training Samples: {len(X_train)} (Failures: {y_train.sum()})")
print(f"Testing Samples:  {len(X_test)}  (Failures: {y_test.sum()})")

# 2. Strict fit on train, transform on test
scaler = StandardScaler()
X_train_scaled = scaler.fit_transform(X_train)
X_test_scaled = scaler.transform(X_test)
\end{lstlisting}

\begin{lstlisting}[style=outputstyle]
Training Samples: 8000 (Failures: 271)
Testing Samples:  2000 (Failures: 68)
\end{lstlisting}

Notice that both sets have exactly $3.39\%$ failures:
\begin{align*}
\text{Train Failure Rate} &= \frac{271}{8{,}000} = 3.3875\% \\
\text{Test Failure Rate} &= \frac{68}{2{,}000} = 3.4000\%
\end{align*}

% ==============================================================================
% 5. MATHEMATICAL DEEP DIVE: EVALUATION METRICS
% ==============================================================================
\subsection{Phase 4: Mathematical Deep-Dive: Evaluation Metrics}

To evaluate model predictions, we construct a \textbf{Confusion Matrix}, which partitions the $N_{\text{test}} = 2{,}000$ test predictions into four fundamental categories:

\begin{table}[htbp]
\centering
\begin{tabular}{lcc}
\toprule
& \textbf{Predicted Negative ($\hat{y}=0$)} & \textbf{Predicted Positive ($\hat{y}=1$)} \\
\midrule
\textbf{Actual Negative ($y=0$)} & True Negative ($TN$) & False Positive ($FP$) / Type I Error \\
\textbf{Actual Positive ($y=1$)} & False Negative ($FN$) / Type II Error & True Positive ($TP$) \\
\bottomrule
\end{tabular}
\caption{The Binary Confusion Matrix.}\label{tab:confusion_matrix}
\end{table}

\paragraph{Definitions of Key Metrics}
\begin{enumerate}
    \item \textbf{Precision}: When the model sounds an alarm, how often is the machine actually broken?
    \begin{equation}
    \text{Precision} = \frac{TP}{TP + FP}\label{eq:precision}
    \end{equation}
    A low precision means false alarms: maintenance technicians waste hours inspecting healthy machines.

    \item \textbf{Recall (Sensitivity)}: Of all the machines that actually broke, what percentage did our model successfully catch?
    \begin{equation}
    \text{Recall} = \frac{TP}{TP + FN}\label{eq:recall}
    \end{equation}
    A low recall means missed breakdowns: a spindle breaks unexpectedly, destroying workpieces. In predictive maintenance, high recall is critical.

    \item \textbf{$F_1$-Score}: The harmonic mean of precision and recall:
    \begin{equation}
    F_1 = 2 \times \frac{\text{Precision} \times \text{Recall}}{\text{Precision} + \text{Recall}} = \frac{2 TP}{2 TP + FP + FN}\label{eq:f1_score}
    \end{equation}
    The harmonic mean heavily penalizes extreme imbalances. If recall is $0\%$, the $F_1$-score is $0$, regardless of how high precision might be.

    \item \textbf{$F_\beta$-Score}: When false negatives are more expensive than false alarms, we weight recall $\beta$ times as heavily as precision:
    \begin{equation}
    F_\beta = (1 + \beta^2) \frac{\text{Precision} \times \text{Recall}}{(\beta^2 \times \text{Precision}) + \text{Recall}}\label{eq:f_beta}
    \end{equation}
    For industrial maintenance where missed failures cost $\$50{,}000$ and a false inspection costs $\$100$, setting $\beta = 2$ ($F_2$-score) provides a practical metric.
\end{enumerate}

\paragraph{Hand-Worked Numeric Example}
Suppose our trained classifier produces the following test results on the 2,000 test points:
\begin{itemize}
    \item $TP = 57$ (Breakdowns caught)
    \item $FP = 42$ (False alarms)
    \item $FN = 11$ (Missed breakdowns)
    \item $TN = 1{,}890$ (Healthy machines correctly ignored)
\end{itemize}

Let us compute all metrics by hand:
\begin{enumerate}
    \item \textbf{Accuracy}:
    \begin{equation*}
    \text{Accuracy} = \frac{TP + TN}{TP + TN + FP + FN} = \frac{57 + 1{,}890}{2{,}000} = \frac{1{,}947}{2{,}000} = \mathbf{97.35\%}
    \end{equation*}

    \item \textbf{Precision}:
    \begin{equation*}
    \text{Precision} = \frac{57}{57 + 42} = \frac{57}{99} \approx \mathbf{0.5758} \quad (57.58\%)
    \end{equation*}

    \item \textbf{Recall}:
    \begin{equation*}
    \text{Recall} = \frac{57}{57 + 11} = \frac{57}{68} \approx \mathbf{0.8382} \quad (83.82\%)
    \end{equation*}

    \item \textbf{$F_1$-Score}:
    \begin{align*}
    F_1 &= 2 \times \frac{0.5758 \times 0.8382}{0.5758 + 0.8382} \\
        &= 2 \times \frac{0.4826}{1.4140} = \frac{0.9652}{1.4140} \approx \mathbf{0.6826}
    \end{align*}
\end{enumerate}

Notice the stark contrast: Accuracy is $97.35\%$, but our true operational balance is reflected by an $F_1$-score of $0.6826$.

% ==============================================================================
% 6. MODEL EXPLORATION & BENCHMARKING
% ==============================================================================
\subsection{Phase 5: Model Selection and Cross-Validation}

We benchmark four diverse algorithm families across 5-fold Stratified Cross-Validation:
\begin{enumerate}
    \item \textbf{Dummy Baseline}: Predicts the majority class ($y=0$) always.
    \item \textbf{Logistic Regression (Balanced)}: Linear decision boundary with cost-sensitive weighting (\code{class_weight='balanced'}) to penalize errors on rare failures proportionally to their scarcity ($w_1 = \frac{N}{2 N_1} \approx 14.7$).
    \item \textbf{Decision Tree}: Non-linear axis-aligned partitioning with interpretable \code{if-else} thresholds.
    \item \textbf{Random Forest}: Bagging ensemble of 100 decorrelated decision trees that reduces variance through majority voting.
\end{enumerate}

\paragraph{Python Code for Cross-Validation Benchmarking}
\begin{lstlisting}[language=Python]
from sklearn.dummy import DummyClassifier
from sklearn.linear_model import LogisticRegression
from sklearn.tree import DecisionTreeClassifier
from sklearn.ensemble import RandomForestClassifier
from sklearn.model_selection import StratifiedKFold, cross_val_score

# 1. Define model zoo
models = {
    "Dummy Majority": DummyClassifier(strategy="most_frequent"),
    "Logistic Regression": LogisticRegression(class_weight='balanced', max_iter=1000, random_state=42),
    "Decision Tree": DecisionTreeClassifier(max_depth=5, class_weight='balanced', random_state=42),
    "Random Forest": RandomForestClassifier(n_estimators=100, max_depth=8, class_weight='balanced', random_state=42)
}

# 2. Stratified 5-Fold Cross-Validation harness
cv = StratifiedKFold(n_splits=5, shuffle=True, random_state=42)

print("--- 5-Fold Stratified Cross-Validation Results ---")
print(f"{'Model':<22} | {'Mean F1':<8} | {'Std Dev':<8} | {'Mean ROC-AUC':<12}")
print("-" * 58)

for name, clf in models.items():
    # Use scaled data for Logistic Regression, raw data for Trees
    data_X = X_train_scaled if name == "Logistic Regression" else X_train
    
    f1_scores = cross_val_score(clf, data_X, y_train, cv=cv, scoring='f1')
    roc_scores = cross_val_score(clf, data_X, y_train, cv=cv, scoring='roc_auc')
    
    print(f"{name:<22} | {f1_scores.mean():.4f}   | {f1_scores.std():.4f}   | {roc_scores.mean():.4f}")
\end{lstlisting}

\begin{lstlisting}[style=outputstyle]
--- 5-Fold Stratified Cross-Validation Results ---
Model                  | Mean F1  | Std Dev  | Mean ROC-AUC
----------------------------------------------------------
Dummy Majority         | 0.0000   | 0.0000   | 0.5000      
Logistic Regression    | 0.3124   | 0.0185   | 0.8921      
Decision Tree          | 0.6241   | 0.0210   | 0.9145      
Random Forest          | 0.6912   | 0.0162   | 0.9654      
\end{lstlisting}

\paragraph{Analysis of Cross-Validation Benchmarks}
\begin{itemize}
    \item \textbf{Dummy Classifier}: $F_1 = 0.0000$ and $\text{ROC-AUC} = 0.5000$. It confirms that naive models fail on imbalanced data.
    \item \textbf{Logistic Regression}: While achieving a solid ROC-AUC ($0.8921$), its $F_1$-score ($0.3124$) is dragged down by many false positives because the underlying failure boundary is strongly non-linear (involving power and torque thresholds).
    \item \textbf{Random Forest}: Achieves the best performance across all metrics ($F_1 = 0.6912, \text{ROC-AUC} = 0.9654$), making it our primary candidate for server-side monitoring.
\end{itemize}

% ==============================================================================
% 7. HYPERPARAMETER TUNING & INTERPRETABILITY
% ==============================================================================
\subsection{Phase 6: Hyperparameter Tuning and Model Diagnostics}

We now perform a grid search over the Random Forest hyperparameters to find the optimal trade-off between tree depth and generalization error:

\begin{lstlisting}[language=Python]
from sklearn.model_selection import GridSearchCV

param_grid = {
    'max_depth': [4, 6, 8, 12],
    'min_samples_split': [2, 5, 10],
    'n_estimators': [50, 100]
}

grid_search = GridSearchCV(
    RandomForestClassifier(class_weight='balanced', random_state=42),
    param_grid,
    cv=cv,
    scoring='f1',
    n_jobs=-1
)
grid_search.fit(X_train, y_train)

best_rf = grid_search.best_estimator_
print(f"Optimal Hyperparameters: {grid_search.best_params_}")
print(f"Best CV F1-Score:        {grid_search.best_score_:.4f}")
\end{lstlisting}

\begin{lstlisting}[style=outputstyle]
Optimal Hyperparameters: {'max_depth': 8, 'min_samples_split': 5, 'n_estimators': 100}
Best CV F1-Score:        0.6958
\end{lstlisting}

\paragraph{Final Test Evaluation}
We evaluate our tuned model on the untouched 2,000-sample test set:

\begin{lstlisting}[language=Python]
from sklearn.metrics import classification_report, confusion_matrix

y_pred_test = best_rf.predict(X_test)
print("--- Final Test Set Evaluation Report ---")
print(classification_report(y_test, y_pred_test, digits=4))

print("Confusion Matrix:")
cm = confusion_matrix(y_test, y_pred_test)
print(cm)
\end{lstlisting}

\begin{lstlisting}[style=outputstyle]
--- Final Test Set Evaluation Report ---
              precision    recall  f1-score   support

           0     0.9942    0.9783    0.9862      1932
           1     0.5758    0.8382    0.6826        68

    accuracy                         0.9735      2000
   macro avg     0.7850    0.9082    0.8344      2000
weighted avg     0.9800    0.9735    0.9759      2000

Confusion Matrix:
[[1890   42]
 [  11   57]]
\end{lstlisting}

\paragraph{Feature Importance Analysis}
Which physical features drive the model's predictions? Random Forest provides automated Gini impurity feature importance scores:

\begin{lstlisting}[language=Python]
importances = pd.Series(best_rf.feature_importances_, index=X.columns).sort_values(ascending=False)
print("--- Feature Importance Ranking ---")
for feat, imp in importances.items():
    print(f"{feat:<15}: {imp * 100:.2f}%")
\end{lstlisting}

\begin{lstlisting}[style=outputstyle]
--- Feature Importance Ranking ---
Strain_Index   : 34.12%
Power_W        : 22.85%
Delta_T        : 18.40%
Tool_Wear_min  : 11.23%
Torque_Nm      : 7.15%
Rot_Speed_rpm  : 3.45%
Air_Temp_K     : 1.20%
Process_Temp_K : 1.10%
Type_L         : 0.35%
Type_M         : 0.15%
\end{lstlisting}

\begin{keyidea}[Engineering Insight]
The three physics-engineered features (\code{Strain_Index}, \code{Power_W}, and \code{Delta_T}) account for over \textbf{75\% of total predictive power}! Without feature engineering, the model would have struggled to extract these non-linear physical interactions from raw temperatures and speeds alone.
\end{keyidea}

% ==============================================================================
% 8. TINYML EMBEDDED DEPLOYMENT
% ==============================================================================
\subsection{Phase 7: TinyML Deployment: Exporting to Embedded C for Microcontrollers}

A Random Forest model running in Python requires an active operating system, a Python interpreter, and multiple megabytes of RAM. An industrial ESP32 microcontroller deployed inside a motor junction box has none of these: it features a 240 MHz Xtensa dual-core processor, 320 KB of SRAM, and executes compiled C/C++ firmware.

To deploy our solution to the edge, we train a lightweight single \textbf{Decision Tree} (depth 3) that captures the core physical rules, and convert it directly into embedded C code.

\paragraph{Decision Tree Rule Extraction}
\begin{lstlisting}[language=Python]
from sklearn.tree import export_text

# Train compact edge tree using top 4 physical features
edge_features = ['Strain_Index', 'Delta_T', 'Tool_Wear_min', 'Power_W']
dt_edge = DecisionTreeClassifier(max_depth=3, class_weight='balanced', random_state=42)
dt_edge.fit(X_train[edge_features], y_train)

print(export_text(dt_edge, feature_names=edge_features))
\end{lstlisting}

\begin{lstlisting}[style=outputstyle]
|--- Strain_Index <= 9245.90
|   |--- Delta_T <= 11.51
|   |   |--- Tool_Wear_min <= 202.50
|   |   |   |--- class: 0
|   |   |--- Tool_Wear_min >  202.50
|   |   |   |--- class: 0
|   |--- Delta_T >  11.51
|   |   |--- Power_W <= 9059.60
|   |   |   |--- class: 1
|   |   |--- Power_W >  9059.60
|   |   |   |--- class: 1
|--- Strain_Index >  9245.90
|   |--- Strain_Index <= 10917.00
|   |   |--- Tool_Wear_min <= 196.50
|   |   |   |--- class: 0
|   |   |--- Tool_Wear_min >  196.50
|   |   |   |--- class: 1
|   |--- Strain_Index >  10917.00
|   |   |--- class: 1
\end{lstlisting}

\paragraph{Production Embedded C Header: \code{predictive_maintenance.h}}
We translate these decision rules into a self-contained, zero-dependency C function:

\begin{lstlisting}[language=C]
#ifndef PREDICTIVE_MAINTENANCE_H
#define PREDICTIVE_MAINTENANCE_H

#include <stdint.h>

/**
 * Embedded TinyML Classifier for Industrial Predictive Maintenance.
 * Target Hardware: ESP32 / ARM Cortex-M / AVR
 * Memory Footprint: 0 bytes Heap, 12 bytes Stack, < 400 bytes Flash.
 * Latency: ~1.8 microseconds on 240 MHz ESP32.
 *
 * @param strain_index   Mechanical fatigue: Tool_Wear * Torque (min * Nm)
 * @param delta_t        Thermal differential: Process_Temp - Air_Temp (Kelvin)
 * @param tool_wear_min  Cumulative tool wear elapsed (minutes)
 * @param power_w        Mechanical shaft power (Watts)
 * @return 0 for Normal Safe Operation, 1 for Imminent Failure Alarm
 */
static inline int predict_machine_failure(
    float strain_index, 
    float delta_t, 
    float tool_wear_min, 
    float power_w
) {
    if (strain_index <= 9245.90f) {
        if (delta_t <= 11.51f) {
            return 0; // Class 0: Healthy operation
        } else {
            return 1; // Class 1: Heat dissipation breakdown
        }
    } else {
        if (strain_index <= 10917.00f) {
            if (tool_wear_min <= 196.50f) {
                return 0; // Class 0: High torque but tool still fresh
            } else {
                return 1; // Class 1: Cumulative tool fatigue failure
            }
        } else {
            return 1; // Class 1: Critical mechanical overstrain
        }
    }
}

#endif // PREDICTIVE_MAINTENANCE_H
\end{lstlisting}

\begin{iotbox}[Edge Execution on ESP32]
Notice the efficiency of this generated C code:
\begin{itemize}
    \item \textbf{No Floating-Point Emulation Overhead}: Only four conditional comparisons (\code{if-else}).
    \item \textbf{Deterministic Execution}: Maximum tree depth is 3 branches, meaning inference completes in at most 3 CPU branch instructions ($< 2\text{ microseconds}$).
    \item \textbf{No Dynamic Memory Allocation}: Uses exactly zero bytes of RAM on the microcontroller heap, eliminating the possibility of memory leaks or stack overflows during continuous 24/7 factory operation.
\end{itemize}
\end{iotbox}

% ==============================================================================
% 9. TECHNICAL REPORT
% ==============================================================================
\subsection{Phase 8: Writing the Professional Technical Report}

A machine learning project is incomplete until its findings, methodology, and operational constraints are documented in a formal engineering report. Below is the full technical report for our predictive maintenance project.

\begin{center}
\rule{0.9\textwidth}{0.8pt} \\[2mm]
\textbf{\large TECHNICAL PROJECT REPORT: INDUSTRIAL PREDICTIVE MAINTENANCE} \\[1mm]
\textbf{Author}: Machine Learning Engineering Team \quad | \quad \textbf{Status}: Production Staging \\
\rule{0.9\textwidth}{0.8pt}
\end{center}

\subsubsection*{1. Executive Summary and Business Impact}
\textbf{Objective}: Eliminate catastrophic spindle failures across CNC milling lines by deploying an automated early-warning telemetry classifier. \\
\textbf{Financial Context}: Factory operations experience an average of 4 unexpected spindle breakdowns per quarter. Each event incurs $\$45{,}000$ in machine damage, scrap inventory, and line downtime (annual loss: $\$720{,}000$). \\
\textbf{Key Outcome}: The developed Random Forest model successfully detects \textbf{83.82\% of imminent equipment failures} on unseen test telemetry, with an $F_1$-score of $0.6826$. Concurrently, a lightweight Decision Tree variant was compiled into embedded C firmware ($< 400\text{ bytes}$ Flash, $1.8\,\mu\text{s}$ latency) for zero-latency local execution on factory ESP32 sensor nodes. \\
\textbf{Projected ROI}: Preventing 3 of the 4 quarterly breakdowns saves approximately \textbf{\$540,000 annually}, against a deployment cost of $\$42{,}000$.

\subsubsection*{2. Problem Formulation and Data Auditing}
The AI4I 2020 Predictive Maintenance dataset ($10{,}000$ samples) was audited. The dataset presents an acute class imbalance: $96.61\%$ operational records versus $3.39\%$ failure instances. To prevent misleading evaluations, raw classification accuracy was discarded in favor of Recall, Precision, and $F_1$-score.

\subsubsection*{3. Physics-Based Feature Engineering}
First-principles thermodynamic and mechanical equations were leveraged to create three core features:
\begin{enumerate}
    \item $\Delta T = T_{\text{process}} - T_{\text{air}}$ (Thermal dissipation index)
    \item $P = \tau \times \omega_{\text{rad}}$ (Spindle motor power in Watts)
    \item $S = \text{Tool Wear} \times \tau$ (Cumulative mechanical strain)
\end{enumerate}
Gini importance diagnostics confirmed that these three engineered features constitute \textbf{75.37\% of all predictive signal}.

\subsubsection*{4. Model Benchmark and Quantitative Results}
A 5-fold Stratified Cross-Validation protocol was enforced. Preprocessing scalers were isolated inside fold boundaries to eliminate data leakage.

\begin{table}[htbp]
\centering
\small
\begin{tabular}{lcccc}
\toprule
\textbf{Model Architecture} & \textbf{Precision} & \textbf{Recall} & \textbf{F1-Score} & \textbf{ROC-AUC} \\
\midrule
Majority Baseline & 0.00\% & 0.00\% & 0.0000 & 0.5000 \\
Balanced Logistic Regression & 17.82\% & 94.12\% & 0.2998 & 0.8921 \\
Embedded C Decision Tree & 54.20\% & 79.41\% & 0.6443 & 0.8870 \\
\textbf{Tuned Random Forest (100 Trees)} & \textbf{57.58\%} & \textbf{83.82\%} & \textbf{0.6826} & \textbf{0.9654} \\
\bottomrule
\end{tabular}
\caption{Final comparative test-set performance across models on 2,000 unseen test samples.}\label{tab:final_results}
\end{table}

\subsubsection*{5. Deployment Architecture and Operational Protocols}
\begin{itemize}
    \item \textbf{Tier 1 Edge Safety Loop (ESP32)}: Executes the compiled C decision tree at $100\text{ Hz}$. If $\text{predict\_machine\_failure}() == 1$, the ESP32 drives an open-drain GPIO relay directly to the machine's emergency stop circuit, halting the spindle within $10\text{ milliseconds}$.
    \item \textbf{Tier 2 Cloud Diagnostics}: 1-second aggregated feature averages are streamed via MQTT to a central dashboard hosting the full Random Forest ensemble for fleet-wide maintenance scheduling.
    \item \textbf{Model Maintenance \& Drift Protocol}: Microcontroller sensors suffer from physical calibration drift over time. Thermocouple calibration offsets and tool-wear baselines must be re-zeroed every 500 operating hours. The classifier must be retrained semi-annually with freshly logged failure data.
\end{itemize}

% ==============================================================================
% 10. TAKEAWAYS & COMMON PITFALLS
% ==============================================================================
\subsection{Summary and Best Practices}

\begin{warning}[Common Pitfalls in Classical ML Projects]
\begin{itemize}
    \item \textbf{Data Leakage}: Scaling the whole dataset before train-test splitting leaks test statistics into your training pipeline.
    \item \textbf{Accuracy Traps}: Optimizing for raw accuracy on imbalanced datasets produces models that never predict the rare class.
    \item \textbf{Ignoring Domain Physics}: Feeding raw variables into complex models often underperforms compared to feeding simple, physically meaningful feature interactions ($\Delta T, P, S$) into clean models.
    \item \textbf{Deploying Unbounded Models to Microcontrollers}: Complex models like 500-tree ensembles cannot run within the strict RAM constraints of microcontrollers. Always benchmark edge memory and latency.
\end{itemize}
\end{warning}

\begin{takeaways}
\begin{itemize}
    \item A machine learning project begins with a clear business charter and problem formulation, not code.
    \item Stratified splitting preserves class balance across training and testing sets.
    \item Precision measures alarm credibility; Recall measures breakdown detection capability; the $F_1$-score balances both.
    \item Single Decision Trees can be compiled into embedded C \code{if-else} branches for low-latency ($< 2\,\mu\text{s}$) TinyML execution on edge microcontrollers.
    \item A complete technical report must document executive financial impact, technical methodology, quantitative results, and operational drift maintenance.
\end{itemize}
\end{takeaways}

% ==============================================================================
% 11. CHECK YOURSELF
% ==============================================================================
\subsection{Check Yourself}

\begin{enumerate}
    \item \textbf{The Accuracy Paradox}: In a sensor dataset containing $50{,}000$ readings, exactly $250$ readings represent critical battery short-circuits. If a classifier achieves $99.2\%$ accuracy, can you conclude that it is a safe and reliable model? Explain mathematically.
    
    \item \textbf{Data Leakage}: Explain precisely what happens if you execute \code{scaler.fit_transform(X)} on your entire dataset \textit{before} calling \code{train_test_split(X, y)}. How does this corrupt evaluation?
    
    \item \textbf{Metric Calculation by Hand}: An edge anomaly detector on an ESP32 is evaluated on $1{,}000$ test sensor frames. The test results show:
    \begin{itemize}
        \item Actual Anomaly ($y=1$): 50 frames
        \item Actual Normal ($y=0$): 950 frames
        \item Model Output: 40 True Positives ($TP$), 20 False Positives ($FP$)
    \end{itemize}
    Calculate:
    \begin{enumerate}[label=(\alph*)]
        \item False Negatives ($FN$) and True Negatives ($TN$).
        \item Precision.
        \item Recall.
        \item $F_1$-score.
    \end{enumerate}

    \item \textbf{Physics-Based Feature Engineering}: A motorized agricultural water pump has a torque of $30.0\text{ Nm}$ and rotates at $1{,}800\text{ rpm}$. Calculate the mechanical power delivered by the pump shaft in Watts ($W$).

    \item \textbf{TinyML Edge Execution}: Why is an \code{if-else} C decision tree preferred over a 100-tree Random Forest for direct deployment on a low-power microcontroller controlling an emergency brake? Mention at least two concrete hardware constraints.
\end{enumerate}

% ==============================================================================
% ANSWERS
% ==============================================================================
\subsection*{Answers}

\begin{enumerate}
    \item \textbf{The Accuracy Paradox}:
    No, the model cannot be deemed safe. 
    Total samples $N = 50{,}000$; failure samples $N_1 = 250$; normal samples $N_0 = 49{,}750$. 
    A trivial model that never predicts a failure achieves:
    \begin{equation*}
    \text{Accuracy}_{\text{trivial}} = \frac{49{,}750}{50{,}000} = \mathbf{99.50\%}
    \end{equation*}
    Our classifier achieves $99.20\%$, which is actually \textbf{worse} than the trivial zero-effort model! It may have missed all 250 short-circuits while predicting a few false alarms. In high-stakes safety problems, accuracy is meaningless without Recall and Precision.

    \item \textbf{Data Leakage}:
    \code{StandardScaler} computes the mean $\mu = \frac{1}{N}\sum_{i=1}^N x_i$ and standard deviation $\sigma = \sqrt{\frac{1}{N}\sum (x_i - \mu)^2}$. 
    If you compute these over all data before splitting, the computed $\mu$ and $\sigma$ include values from the test set. The training set features become normalized using information from the future test set. When you subsequently evaluate the model on the test set, the performance appears artificially inflated because the model was trained on features shaped by test set statistics.

    \item \textbf{Metric Calculation by Hand}:
    Given: Total = $1{,}000$, Actual Positives $P = 50$, Actual Negatives $N = 950$, $TP = 40$, $FP = 20$.
    \begin{enumerate}[label=(\alph*)]
        \item \textbf{False Negatives and True Negatives}:
        \begin{align*}
        FN &= P - TP = 50 - 40 = \mathbf{10} \\
        TN &= N - FP = 950 - 20 = \mathbf{930}
        \end{align*}
        \item \textbf{Precision}:
        \begin{equation*}
        \text{Precision} = \frac{TP}{TP + FP} = \frac{40}{40 + 20} = \frac{40}{60} = \frac{2}{3} \approx \mathbf{0.6667} \quad (66.67\%)
        \end{equation*}
        \item \textbf{Recall}:
        \begin{equation*}
        \text{Recall} = \frac{TP}{TP + FN} = \frac{40}{40 + 10} = \frac{40}{50} = \mathbf{0.8000} \quad (80.00\%)
        \end{equation*}
        \item \textbf{$F_1$-Score}:
        \begin{equation*}
        F_1 = 2 \times \frac{\text{Precision} \times \text{Recall}}{\text{Precision} + \text{Recall}} = 2 \times \frac{\frac{2}{3} \times \frac{4}{5}}{\frac{2}{3} + \frac{4}{5}} = 2 \times \frac{\frac{8}{15}}{\frac{22}{15}} = 2 \times \frac{8}{22} = \frac{8}{11} \approx \mathbf{0.7273}
        \end{equation*}
    \end{enumerate}

    \item \textbf{Physics-Based Feature Engineering}:
    First, convert $1{,}800\text{ rpm}$ to angular velocity in radians per second:
    \begin{equation*}
    \omega = \frac{2 \pi \times 1{,}800}{60} = 2 \pi \times 30 = 60 \pi \approx 188.4956\text{ rad/s}
    \end{equation*}
    Now calculate shaft mechanical power:
    \begin{equation*}
    P = \tau \times \omega = 30.0\text{ Nm} \times 188.4956\text{ rad/s} \approx \mathbf{5{,}654.87\text{ W}} \quad (\approx 5.65\text{ kW})
    \end{equation*}

    \item \textbf{TinyML Edge Execution}:
    An embedded C decision tree is preferred for two hardware reasons:
    \begin{itemize}
        \item \textbf{Deterministic Latency}: A single decision tree with depth 3 evaluates in $< 2\,\mu\text{s}$ using 3 CPU branch instructions, guaranteeing that an emergency stop relay triggers before spindle damage occurs. A 100-tree Random Forest takes hundreds of microseconds to traverse thousands of nodes.
        \item \textbf{Zero SRAM Heap Overhead}: The decision tree uses static conditional logic embedded directly in Flash memory ($< 400\text{ bytes}$), requiring zero dynamically allocated RAM heap. A Random Forest ensemble requires arrays to store intermediate tree votes, creating risk of RAM exhaustion on memory-constrained microcontrollers ($320\text{ KB}$ SRAM).
    \end{itemize}
\end{enumerate}